% !TeX root = ../main.tex
\section{Symbolic Semantics}
\label{sec:proof}

\todo[inline]{Add Lean references}

To motivate our proof strategy, first observe that at an $\E$-mode distribution, the original program samples a value whereas the determinized program computes the mean of the distribution. The resulting programs therefore have different output distributions. Thus, our proof relates both programs to a common symbolic semantics, after which two interpretations recover the source and determinized programs. Our proof thus proceeds in three stages:

\begin{enumerate}
  \item Define the symbolic semantics;
  \item Define the actual and expected interepretations that recover the source and determinized programs, respectively, and prove properties regarding these interpretations; and
  \item Derive the tracewise and global soundness results.
\end{enumerate}

\label{sec:symbolic-semantics}

We first introduce a symbolic semantics that records $\E$-mode
samples without choosing their values, while executing
$\G$-mode samples normally. This provides a common execution
model for the source and determinized programs.

\begin{definition}[Symbolic state]
A symbolic state has the form
\begin{align}
  a
  &=c_0+\sum_{i=1}^{k}c_i u_i,
  \qquad c_0,\ldots,c_k\in\Real,\\
  \sigma
  &::=\emptyset\mid\sigma,\ u\sim D(\bar a),\\
  s
  &::=(\sigma\mid\widehat e).
\end{align}
where $a$ is an affine expression over symbolic variables
$u_1,\ldots,u_k$, $\sigma$ is an ordered symbolic environment,
and $\widehat e$ is a residual symbolic expression. Each declaration $u\sim D(\bar a)$ introduces
a fresh symbolic variable whose distribution parameters
refer only to variables declared earlier in $\sigma$.
\end{definition}

In particular, residual symbolic expressions have the same constructors as source expressions,
with real literals replaced by $\sym a$, where $a$ is an affine expression, to represent
computations involving samples whose values have not yet been chosen.
For example, $\sym2*\sym{u}+\sym1$ can reduce to $\sym{2u+1}$
without sampling $u$; its declaration in $\sigma$ retains the
distribution for later interpretation using sampled values or
sequential means. Typing rules are unchanged except that every
$\sym a$ may have type $\tyR\E$, while only $\sym a$ with constant
$a$ may have type $\tyR\G$. The full syntax and
typing rules appear in \cref{sec:appendix-symbolic-definitions}.

\paragraph{Lifting, realization, and initialization.}
We connect source expressions and residual expressions through two
structural maps, whose full definitions are given in
\cref{def:appendix-symbolic-maps}:

\begin{enumerate}
    \item \emph{Lifting}.
    $\mathsf{lift}:Expr\to\widehat{Expr}$
    embeds a source expression into the residual language
    by replacing each real literal with a constant affine expression:
    \[
        \mathsf{lift}(r)=\sym{r+\sum_i0u_i}.
    \]

    \item \emph{Realization}.
    $\mathsf{realize}_\rho:\widehat{Expr}\to Expr$
    maps a residual expression back to a source expression
    by evaluating each symbolic literal under a valuation $\rho$:
    \[
        \mathsf{realize}_\rho
        \left(\sym{c_0+\sum_{i=1}^k c_i u_i}\right)
        =c_0+\sum_{i=1}^k c_i\rho(u_i).
    \]
\end{enumerate}

Both maps otherwise preserve constructors and distribution modes,
acting recursively on subexpressions. Now, from here, symbolic execution is initialized with the lifted source expression
and an empty symbolic environment:
\[
    \mathsf{init}(e):=(\emptyset\mid\mathsf{lift}(e)).
\]
Since lifting introduces no symbolic variables, realizing the
initial expression under the empty valuation $()$ recovers
the original source expression, $\mathsf{realize}_{()}(\mathsf{lift}(e))=e$.

\begin{figure*}[t]
\centering
\small
\setlength{\jot}{3pt}

\begin{tabular}{@{}p{0.36\textwidth}!{\color{black!20}\vrule width 0.4pt}p{0.58\textwidth}@{}}
\begin{minipage}[t]{\linewidth}
\vspace{0pt}
\centering
\textbf{Residual expression reduction}
\[
\begin{aligned}
  (\lambda x.\widehat e)\,\widehat v
  &\rightsquigarrow_{\mathrm{sym}}
  \widehat e[x\mapsto\widehat v],
  \\[2pt]
  \letin x{\widehat v}{\widehat e}
  &\rightsquigarrow_{\mathrm{sym}}
  \widehat e[x\mapsto\widehat v],
  \\[2pt]
  \avg D{\bar a}
  &\rightsquigarrow_{\mathrm{sym}}
  \operatorname{Mean}^{\#}_D(\bar a).
\end{aligned}
\]
\raggedright
The remaining reduction
rules are lifted from \cref{fig:semantics}. The affine expression
$\operatorname{Mean}^{\#}_D(\bar a)$ represents the primitive mean at the
symbolic arguments $\bar a$.
\end{minipage}
&
\begin{minipage}[t]{\linewidth}
\vspace{0pt}
\centering
\textbf{Symbolic small-step semantics}
\begin{mathpar}
  \inferrule*[right=Reduce]
    {\widehat e_1\rightsquigarrow_{\mathrm{sym}}\widehat e_2}
    {
      (\sigma\mid\widehat C[\widehat e_1]:\rho)
      \xrightarrow{}_{\mathrm{sym}}
      (\sigma\mid\widehat C[\widehat e_2]:\rho)
    }
  \and
  \inferrule*[right=Sample-E]
    {u\notin\operatorname{dom}(\sigma)}
    {
      (\sigma\mid\widehat C[\draw D\E{\bar a}]:\rho)
      \xrightarrow{}_{\mathrm{sym}}
      (\sigma,u\sim D(\bar a)\mid\widehat C[u]:\rho)
    }
  \and
  \inferrule*[right=Sample-G]
    {\operatorname{Domain}_D(\bar r)}
    {
      (\sigma\mid\widehat C[\draw D\G{\bar r}]:\rho)
      \xrightarrow{}_{\mathrm{sym}}
      (\sigma\mid\widehat C[r]:\rho)
    }
\end{mathpar}
\raggedright
The \textsc{Reduce} rule lifts reduction of residual symbolic expressions
to reduction of symbolic states under a symbolic evaluation context.
In \textsc{Sample-G}, the label $r$ is distributed according to
$D(\bar r)$. The label $\epsilon$ denotes a step that contributes no entry
to the $\G$ trace.
\end{minipage}
\end{tabular}

\caption{Symbolic semantics. Symbolic expression reduction
$\rightsquigarrow_{\mathrm{sym}}$ on residual expressions $\widehat e$ is
shown on the left, and the corresponding small-step transition relation
$\xrightarrow{}_{\mathrm{sym}}$ on symbolic states is shown on the right.
The \textsc{Reduce} rule lifts residual expression reduction through a
symbolic evaluation context, while \textsc{Sample-E} and \textsc{Sample-G}
handle the two sampling modes directly. Symbolic evaluation contexts
$\widehat C$ have the same grammar as the evaluation contexts $C$ in
\cref{fig:semantics}.}
\label{fig:symbolic-semantics}
\end{figure*}

\Cref{fig:symbolic-semantics} presents selected residual
reduction rules and the symbolic state transitions.
The full residual reduction and symbolic transition rules
appear in \cref{sec:appendix-symbolic-rules}. Deterministic reductions follow the concrete semantics,
performing arithmetic on affine expressions.
At an $\E$-mode expression $\draw D\E{\sym{\bar a}}$, symbolic
execution introduces a fresh variable $u$, appends
$u\sim D(\bar a)$ to $\sigma$, and replaces the expression
by $\sym u$. This records the draw without sampling it or
replacing it by its mean. A $\G$-mode distribution is instead sampled immediately.

\subsection{Two Interpretations of the Symbolic Semantics}
\label{sec:symbolic-interpretations}

A symbolic state admits two interpretations: the
\emph{actual interpretation} samples the delayed $\E$-mode variables,
whereas the \emph{expected interpretation} replaces them by their
sequential means.

\begin{definition}[Actual and expected interpretations]
Let
\[
\begin{aligned}
  s
  &=
  (\sigma\mid\widehat e:\rho),
  \qquad
  \sigma
  =
  \bigl(
    u_1\sim D_1(\bar a_1),\ldots,
    u_k\sim D_k(\bar a_k)
  \bigr),
  \\
  \rho\sim\mathcal A_\sigma
  &:\qquad
  \rho(u_i)
  \sim
  D_i\bigl(\bar a_i[\rho]\bigr),
  \qquad i=1,\ldots,k
  \quad\text{in declaration order},
  \\
  \bar\rho_\sigma(u_i)
  &:=
  \operatorname{Mean}_{D_i}
  \bigl(\bar a_i[\bar\rho_\sigma]\bigr),
  \qquad i=1,\ldots,k
  \quad\text{in declaration order},
  \\
  \mathsf{Act}_n(s)
  &:=
  \int
    \joint{\mathsf{realize}_\rho(\widehat e)}n\,
    \mathcal A_\sigma(d\rho),
  \\
  \mathsf{Exp}_n(s)
  &:=
  \joint{
    \trans{
      \mathsf{realize}_{\bar\rho_\sigma}(\widehat e)
    }
  }n.
\end{aligned}
\]
Here $\mathcal A_\sigma$ is the joint distribution of the valuation $\rho$ obtained
by sampling the declarations in $\sigma$ from left to right. Because
$\bar a_i$ depends only on variables declared before $u_i$, both
$\mathcal A_\sigma$ and the sequential mean valuation
$\bar\rho_\sigma$ are well defined in declaration order.
\end{definition}

Both $\mathsf{Act}_n(s)$ and $\mathsf{Exp}_n(s)$ are measures on
$\Trace\times\Real$. The first recovers the source behavior, while the
second recovers the determinized behavior. The proof shows that each symbolic step corresponds to a step in both
programs under their respective interpretations. For a $\G$-mode sample,
both programs sample from the same distribution. For an $\E$-mode sample,
the symbolic step records a new variable: the actual interpretation
samples its value, while the expected interpretation uses its sequential
mean. Thus, one symbolic execution relates the original and determinized
programs step by step. Traces are recorded by the interpretations,
rather than stored in symbolic states.

\begin{lemma}[Interpretations commute with a symbolic step]
\label{lem:interpretation-commuting}
Let $s$ be a well-typed, domain-safe symbolic state. If $s\xrightarrow{}_{\mathrm{sym}}s'$, then, for every $n\in\Nat$,
\[
  \mathsf{Act}_{n+1}(s)=\mathsf{Act}_n(s'),
  \qquad
  \mathsf{Exp}_{n+1}(s)=\mathsf{Exp}_n(s').
\]
For a $\G$-mode state $s=
  \bigl(
    \sigma
    \mid
    \widehat C[\draw D\G{\sym{\bar r}}]
    :\rho
  \bigr)$,
every measurable $A\subseteq\Trace$ and $B\subseteq\Real$ satisfy
\[
\begin{aligned}
\mathsf{Act}_{n+1}(s)(A\times B)
&=
\int_{\Real}
  \mathsf{Act}_n
  \bigl(\sigma\mid\widehat C[\sym r]:\rho\bigr)
  \bigl(
    \{(\tau,x)\mid (D,r)::\tau\in A,\ x\in B\}
  \bigr)
  \,D(\bar r)(dr),
\\
\mathsf{Exp}_{n+1}(s)(A\times B)
&=
\int_{\Real}
  \mathsf{Exp}_n
  \bigl(\sigma\mid\widehat C[\sym r]:\rho\bigr)
  \bigl(
    \{(\tau,x)\mid (D,r)::\tau\in A,\ x\in B\}
  \bigr)
  \,D(\bar r)(dr).
\end{aligned}
\]
\end{lemma}

Initializing with an empty environment and using
$\mathsf{realize}_{()}(\mathsf{lift}(e))=e$ connects the two
interpretations to the concrete semantics:

\begin{theorem}[Interpretation of initialization]
\label{thm:interpretation-initialization}
Suppose $\vdash e:\tyR\E$ and $e$ is closed. Then, for every $n\in\Nat$,
\[
  \mathsf{Act}_n\bigl(\mathsf{init}(e)\bigr)=\joint e n,
  \qquad
  \mathsf{Exp}_n\bigl(\mathsf{init}(e)\bigr)=\joint{\trans e}n.
\]
\end{theorem}


\begin{theorem}[Agreement of the two interpretations]
\label{thm:interpretation-agreement}
Let $s$ be a well-typed, domain-safe symbolic state. For every
$n\in\Nat$, $\mathsf{Act}_n(s)$ and $\mathsf{Exp}_n(s)$ have the same
trace marginal: for every measurable $A\subseteq\Trace$,
\[
  \mathsf{Act}_n(s)(A\times\Real)
  =
  \mathsf{Exp}_n(s)(A\times\Real).
\]
Moreover, for
$\mathsf{Act}_n(s)(\,\cdot\times\Real)$-almost every trace $\tau$,
\[
  \int_{\Real}
    |r|\,\mathsf{Act}_n(s)(dr\mid\tau)
  <\infty,
  \qquad
  \mathsf{Exp}_n(s)(\,\cdot\mid\tau)
  =
  \delta_{
    \displaystyle
    \int_{\Real}
      r\,\mathsf{Act}_n(s)(dr\mid\tau)
  }.
\]
\end{theorem}

\Cref{fig:symbolic-example} works through these results on the program
$e=\draw{uniform}\E{1,2}\,/\,\draw{uniform}\G{0,1}$. The symbolic
execution delays the $\E$-draw as a variable $u_1$, samples the $\G$-draw
as $r$, and ends in the residual value $\sym{\tfrac1r u_1}$. Realizing $u_1$ as a
sample $c$ recovers each step of the source program, and realizing it as
its mean $\tfrac32$ recovers each step of $\trans e$. For each fixed trace
$[(\kw{uniform},r)]$, the conditional mean $\tfrac3{2r}$ is finite.
Integrating over $r$, however, gives a global expectation of $+\infty$
for both programs, which is consistent with \cref{thm:expectation}.

% !TeX root = ../main.tex
\begin{figure}[t]
\centering
\small
\newcommand{\U}{\kw{unif}}
\let\dist\draw
\newcommand{\tr}[1]{{\color{black!55}#1}}
\newcommand{\nil}{\tr{[\,]}}
\newcommand{\tg}{\tr{[(\U,r)]}}
\begin{tikzpicture}[
  every node/.style={font=\small},
  col/.style={inner sep=2pt},
  hdr/.style={align=center, font=\footnotesize, inner sep=1pt},
]
\matrix (m) [matrix of math nodes, row sep=2.5mm, column sep=4.5mm,
             nodes={col, anchor=center}] {
  (\nil,\ \dist\U\E{1,2}\,/\,\dist\U\G{0,1})
  & \bigl(\emptyset\mid\dist\U\E{1,2}\,/\,\dist\U\G{0,1}\bigr)
  & (\nil,\ \dist\U\M{1,2}\,/\,\dist\U\G{0,1}) \\
  (\nil,\ c\,/\,\dist\U\G{0,1})
  & \bigl(u_1\sim\U(1,2)\mid u_1\,/\,\dist\U\G{0,1}\bigr)
  & (\nil,\ \tfrac32\,/\,\dist\U\G{0,1}) \\
  (\tg,\ c\,/\,r)
  & \bigl(u_1\sim\U(1,2)\mid u_1\,/\,r\bigr)
  & (\tg,\ \tfrac32\,/\,r) \\
  (\tg,\ \tfrac{c}{r})
  & \bigl(u_1\sim\U(1,2)\mid \tfrac1r\,u_1\bigr)
  & (\tg,\ \tfrac{3}{2r}) \\
};
% column headers, aligned on a common baseline
\path (m-1-2.north) ++(0,5mm) coordinate (hb);
\node[hdr, anchor=south] (h1) at (hb -| m-1-1)
  {\textbf{Actual} \ $\mathsf{Act}_{k}(s_k)$\\$\mathsf{realize}_\rho(\widehat e_k)$\\$\rho(u_1)=c$};
\node[hdr, anchor=south] (h2) at (hb -| m-1-2)
  {\textbf{Symbolic state}\\$s_k=(\sigma_k\mid\widehat e_k)$\\$s_0=\mathsf{init}(e)$};
\node[hdr, anchor=south] (h3) at (hb -| m-1-3)
  {\textbf{Expected} \ $\mathsf{Exp}_{k}(s_k)$\\$\trans{\mathsf{realize}_{\bar\rho}(\widehat e_k)}$\\$\bar\rho(u_1)=\tfrac32$};
\draw[black!25] ($(hb -| m.west)+(0,-1.5mm)$) -- ($(hb -| m.east)+(0,-1.5mm)$);
% row labels on the symbolic states
\foreach \k in {1,...,4} {
  \pgfmathtruncatemacro\kk{\k-1}
  \node[anchor=east, font=\footnotesize, black!60, inner sep=1pt]
    at (m-\k-2.west) {$s_{\kk}$};
}
% result row: the joint measures read off at s_3, per column
\path (m-4-2.south) ++(0,-4mm) coordinate (rb);
\draw[black!25] (rb -| m.west) -- (rb -| m.east);
\path (rb) ++(0,-1.5mm) coordinate (rt);
\node[hdr, anchor=north] at (rt -| m-4-1)
  {$\mathsf{Act}_3(s_0)=J_e$ \ \textcolor{black!55}{(\cref{thm:interpretation-initialization})}\\[1pt]
   $\law e(\,\cdot\mid\tau_r)=\U\bigl(\tfrac1r,\tfrac2r\bigr)$\\[1pt]
   $\Expect[e\mid\tau_r]=\tfrac{3}{2r}<\infty$};
\node[hdr, anchor=north] at (rt -| m-4-2)
  {$\tau_r:=[(\U,r)]$, \ $r\sim\U(0,1)$\\[1pt]
   $\trlaw e=\trlaw{\trans e}$ \ \textcolor{black!55}{(\cref{lem:trace-preservation})}\\[1pt]
   $\Expect_{\mathrm{ret}}=\int_0^1\tfrac{3}{2r}\,dr=+\infty$ \ \textcolor{black!55}{(\cref{thm:expectation})}};
\node[hdr, anchor=north, text=teal!75!black] at (rt -| m-4-3)
  {$\mathsf{Exp}_3(s_0)=J_{\trans e}$ \ \textcolor{black!55}{(\cref{thm:interpretation-initialization})}\\[1pt]
   $\law{\trans e}(\,\cdot\mid\tau_r)=\delta_{3/(2r)}$\\[1pt]
   $=\delta_{\Expect[e\mid\tau_r]}$ \ \textcolor{black!55}{(\cref{thm:interpretation-agreement})}};
\end{tikzpicture}

\caption{Symbolic execution of
$e=\draw{uniform}\E{1,2}\,/\,\draw{uniform}\G{0,1}$, shown in lockstep with
the source program (left) and $\trans e$ (right). 
% Row $k$ is the $k$-th
% symbolic state $s_k$, and the outer columns are its two realizations:
% $\mathsf{Act}_{3-k}(s_k)$ is the joint measure $\joint{\cdot}{3-k}$ of the
% left entry averaged over $c\sim\kw{uniform}(1,2)$, and
% $\mathsf{Exp}_{3-k}(s_k)$ is $\joint{\cdot}{3-k}$ of the right entry. Each outer entry
% pairs the trace prefix emitted so far (gray) with the current expression,
% matching $J$ on $\Trace\times\Real$; only \textsc{Sample-G} extends the
% trace.
% The bottom row reads off the resulting joint measures for the trace
% $\tau_r$.
% \textsc{Sample-E} delays the $\E$-draw as $u_1$ rather than sampling it
% (left) or averaging it (right). \textsc{Sample-G} replaces the $\G$-draw by a value
% $r$; both interpretations integrate $r$ against $\kw{uniform}(0,1)$ and
% prepend the same trace entry $(\kw{uniform},r)$. Each step is one instance of
% \cref{lem:interpretation-commuting}, so $\mathsf{Act}_3(s_0)=J_e$ and
% $\mathsf{Exp}_3(s_0)=J_{\trans e}$
% (\cref{thm:interpretation-initialization}). At $s_3$, the result is affine
% in $u_1$, so the expected interpretation returns the conditional mean of
% the actual one (\cref{thm:interpretation-agreement}). We abbreviate $\kw{uniform}$ as $\kw{unif}$.
}
\label{fig:symbolic-example}
\Description{A three-column table steps the program uniform E of 1 and 2,
divided by uniform G of 0 and 1, through four states. The middle column is
the symbolic state: the E draw becomes a variable u1 declared with
distribution uniform of 1 and 2, the G draw is sampled as r and recorded in
the trace, and the result is the affine expression u1 over r. The left column
realizes u1 as a sample c, giving c over r; the right column realizes u1 as
its mean 3/2, giving 3 over 2r. Below, the joint measures show conditional
law uniform of 1 over r and 2 over r versus a point mass at 3 over 2r, with
the same trace law and an infinite global expectation on both sides.}
\end{figure}


\subsection{Deriving the Soundness Results}

The commuting argument establishes symbolic soundness for every
$n\in\Nat$. By \cref{thm:interpretation-initialization}, initializing the
symbolic semantics recovers $\joint e n$ under the actual interpretation
and $\joint{\trans e}n$ under the expected interpretation. The proof then
shows that these results use a common conditional output distribution indexed by the
retained trace. Summing over $n$ gives
\[
  J_{\trans e}(A\times B)
  =
  \int_A
    \delta_{\Expect[e\mid\tau]}(B)\,
    \trlaw e(d\tau)
\]
for all measurable $A\subseteq\Trace$ and $B\subseteq\Real$. By \cref{lem:trace-preservation},
$\trlaw{\trans e}=\trlaw e$. The definition of the conditional output distribution
also gives
\[
  J_{\trans e}(A\times B)
  =
  \int_A
    \law{\trans e}(B\mid\tau)\,
    \trlaw e(d\tau).
\]
The two expressions therefore disintegrate the same joint measure with
respect to the same trace measure. Uniqueness of disintegration implies
\[
  \law{\trans e}(\,\cdot\mid\tau)
  =
  \delta_{\Expect[e\mid\tau]}
\]
for $\trlaw e$-almost every trace $\tau$, proving
\cref{thm:trace-soundness}.

Finally, integrating this tracewise characterization over the common trace
measure gives the global results. It yields expectation and
successful return probability preservation, while Jensen's inequality
gives the convex function and variance guarantees.
